\reading{CR-10}{Credit Value at Risk}{STRIKE FIRST}{6}

\begin{crkeybox}[Learning objectives for this reading]
\begin{enumerate}[label=\textbf{\alph*.},leftmargin=2em]
\item Compare market risk value at risk (VaR) with credit VaR in terms of definition, time horizon, and tools for measuring them.
\item Define and calculate Credit VaR.
\item Describe the use of rating transition matrices for calculating credit VaR.
\item Describe the application of the Vasicek's model to estimate capital requirements under the Basel II internal-ratings-based (IRB) approach.
\item Explain the Vasicek's model, Credit Risk Plus (CreditRisk+) model, and the CreditMetrics ways of estimating the probability distribution of losses arising from defaults as well as modeling the default correlation.
\item Define credit spread risk and assess its impact on calculating credit VaR.
\end{enumerate}
\end{crkeybox}

\begin{crtrapbox}[Single-layer reading]
CR-10 exists only in the Crash layer. There is no Class Lecture treatment.
\end{crtrapbox}

% =====================================================================
\lo{CR-10 a}{Compare market risk value at risk (VaR) with credit VaR in terms of definition, time horizon, and tools for measuring them.}

Credit risk VaR is similar to market VaR, but focuses on potential losses from
loan defaults, downgrades, and credit spread changes.

\begin{center}
\footnotesize
\begin{tabular}{L{3.0cm} L{5.6cm} L{5.9cm}}
\toprule
 & \thead{Market VaR} & \thead{Credit VaR} \\
\midrule
Loss driver & Market price movements &
  Loan defaults, downgrades, and credit spread changes \\
\term{Time horizon} & One \textbf{day} & One \textbf{year} \\
\term{Calculation tools} & Historical simulation is workable &
  Requires more complex models than the historical simulation used for market
  VaR \\
\bottomrule
\end{tabular}
\end{center}
\figcap{The horizon difference is a factor of roughly 250 trading days, and it is
the reason the tools diverge: a one-day window has enough usable history for
simple simulation, and a one-year window does not.}

% =====================================================================
\lo{CR-10 b}{Define and calculate Credit VaR.}

Two factors drive the calculation of credit VaR:

\begin{enumerate}[label=\alph*)]
\item \term{Credit correlation.} How defaults of different companies are
interconnected.
\item \term{Rating transitions.} Transition matrices show historical likelihoods
of companies moving between credit ratings, for example AAA to AA.
\end{enumerate}

The worked calculation of credit VaR is carried out at \textbf{CR-11 c}, where
both the perfectly correlated and the independent cases are computed on the same
portfolio.

% =====================================================================
\lo{CR-10 c}{Describe the use of rating transition matrices for calculating credit VaR.}

Rating transition matrices predict credit rating changes over time, and are used
for calculating credit VaR.

\begin{center}
\footnotesize
\begin{tabular}{l r r r r r r r r}
\toprule
\thead{Initial} & \multicolumn{8}{c}{\thead{Year-end rating}} \\
\cmidrule(l){2-9}
\thead{rating} & \thead{AAA} & \thead{AA} & \thead{A} & \thead{BBB} &
\thead{BB} & \thead{B} & \thead{CCC/C} & \thead{D} \\
\midrule
AAA & 89.85 & 9.35 & 0.55 & 0.05 & 0.11 & 0.03 & 0.05 & 0.00 \\
AA & 0.50 & 90.76 & 8.08 & 0.49 & 0.05 & 0.06 & 0.02 & 0.02 \\
A & 0.03 & 1.67 & 92.61 & 5.23 & 0.27 & 0.12 & 0.02 & 0.05 \\
BBB & 0.00 & 0.10 & 3.45 & 91.93 & 3.78 & 0.46 & 0.11 & 0.17 \\
BB & 0.01 & 0.03 & 0.12 & 5.03 & 86.00 & 7.51 & 0.61 & 0.70 \\
B & 0.00 & 0.02 & 0.08 & 0.17 & 5.18 & 85.09 & 5.66 & 3.81 \\
CCC/C & 0.00 & 0.00 & 0.12 & 0.20 & 0.65 & 14.72 & 50.90 & 33.42 \\
D & 0.00 & 0.00 & 0.00 & 0.00 & 0.00 & 0.00 & 0.00 & 100.00 \\
\bottomrule
\end{tabular}
\end{center}
\figcap{One-year rating transition matrix for 1981--2020, in \%. Each row sums to
100 (to rounding). Two structural features are worth noting: the diagonal
dominates everywhere except CCC/C, and the D row is absorbing --- once in
default, the probability of staying there is 100\%.}

\subsection*{Powering the matrix to other horizons}

Matrices can be \term{powered} to project multi-year transitions, or
\term{rooted} for periods shorter than a year.

\begin{crexbox}[Extending and shortening the horizon]
The one-year probability of an AAA-rated issuer remaining AAA is 89.85\%. What is
the probability over three years, and over three months?
\tcblower
\textbf{Expanding.} For three years, raise to the third power:
\[ 0.8985^{3} = 0.7254 \text{, or } 72.54\% \]

\textbf{Shortening.} For three months, take the fourth root:
\[ 0.8985^{1/4} = 0.9736 \text{, or } 97.36\% \]
\end{crexbox}

% =====================================================================
\lo{CR-10 d}{Describe the application of the Vasicek's model to estimate capital requirements under the Basel II internal-ratings-based (IRB) approach.}

Vasicek's Gaussian copula model is a method used on a portfolio of loans to
calculate high percentiles of the distribution of the default rate. It gives the
\term{worst case default rate (WCDR)}.

\begin{crfmlbox}[Worst case default rate, and the capital requirement]
\[ \mathrm{WCDR}(T,X) \;=\; N\!\left(\frac{N^{-1}(\mathrm{PD})
   + \sqrt{\rho}\,N^{-1}(X)}{\sqrt{1-\rho}}\right) \]
\[ \text{Capital requirement} \;=\; \sum_{i=1}^{n}
   \mathrm{WCDR}_i(T,X)\times \mathrm{EAD}_i \times \mathrm{LGD}_i \]
\vk{$T$ = the time horizon; $X$ = the confidence level; PD = probability of
default; $\rho$ = the correlation factor; $N^{-1}$ = the inverse cumulative
normal. Note that $\rho$ appears twice, as $\sqrt{\rho}$ in the numerator and
$\sqrt{1-\rho}$ in the denominator, so raising it pushes the ratio up from both
directions.}
\end{crfmlbox}

\begin{itemize}
\item $\rho$, the correlation factor, should reflect the \term{asset return
correlation} between companies. Return on investment or return on equity can be
used as proxies when direct correlation is difficult to assess.
\item One challenge of Vasicek's model is the \term{lack of incorporation of tail
correlation}, meaning extreme events that fall outside the expected probability
distribution.
\end{itemize}

\begin{center}
\begin{tikzpicture}
\begin{axis}[
  width=12.2cm, height=5.8cm,
  xlabel={Correlation factor $\rho$},
  ylabel={Worst case default rate (\%)},
  xlabel style={font=\sffamily\scriptsize}, ylabel style={font=\sffamily\scriptsize},
  tick label style={font=\sffamily\scriptsize},
  xmin=0, xmax=0.4, ymin=0, ymax=36,
  xtick={0,0.1,0.2,0.3,0.4}, xticklabels={0,0.10,0.20,0.30,0.40},
  grid=major, grid style={FRMRule!50, line width=0.3pt}]
\addplot[draw=FRMNavy, line width=1.2pt, smooth]
  coordinates {(0.000,1.000) (0.020,2.816) (0.040,4.062) (0.060,5.275)
  (0.080,6.500) (0.100,7.750) (0.120,9.033) (0.140,10.352) (0.160,11.711)
  (0.180,13.111) (0.200,14.553) (0.220,16.038) (0.240,17.568) (0.260,19.144)
  (0.280,20.767) (0.300,22.438) (0.320,24.158) (0.340,25.929) (0.360,27.751)
  (0.380,29.626) (0.400,31.556)};
\addplot[draw=FRMGreen, line width=1pt, dashed, sharp plot]
  coordinates {(0,1) (0.4,1)};
\node[font=\sffamily\scriptsize\bfseries, text=FRMGreen, fill=white, inner sep=1.5pt]
  at (axis cs:0.075,3.1) {PD $=1\%$};
\addplot[only marks, mark=*, mark size=2pt, draw=FRMRed, fill=FRMRed]
  coordinates {(0.20,14.553)};
\node[font=\sffamily\scriptsize\bfseries, text=FRMRed, fill=white, inner sep=1.5pt]
  at (axis cs:0.267,11.2) {$\rho=0.20 \Rightarrow 14.55\%$};
\draw[-{Stealth[length=4pt]}, FRMRed] (axis cs:0.238,12.1) -- (axis cs:0.205,14.0);
\end{axis}
\end{tikzpicture}
\end{center}
\figcap{WCDR at a 99.9\% confidence level for a portfolio with a 1\% probability
of default, plotted against the correlation factor. At zero correlation the worst
case equals the PD itself. A correlation of only 0.20 --- a modest number ---
multiplies the worst case default rate by more than fourteen. This is why the
capital requirement is so sensitive to $\rho$, and why the model's inability to
capture tail correlation matters.}

% =====================================================================
\lo{CR-10 e}{Explain the Vasicek's model, Credit Risk Plus (CreditRisk+) model, and the CreditMetrics ways of estimating the probability distribution of losses arising from defaults as well as modeling the default correlation.}

\begin{center}
\footnotesize
\begin{tabular}{L{2.6cm} L{5.6cm} L{6.3cm}}
\toprule
\thead{Model} & \thead{What distribution it estimates} &
\thead{How it handles correlation} \\
\midrule
\term{Vasicek} & The probability distribution of \emph{losses} &
  A single systematic factor influencing default probabilities \\
\term{CreditRisk+} & The probability distribution of the \emph{number of
  defaults}: Poisson, or negative binomial &
  Through uncertainty in the default rate itself, rather than an explicit
  correlation parameter \\
\term{CreditMetrics} & The probability distribution of losses, using a
  \emph{ratings transition matrix}. Considers the impact of \textbf{downgrades}
  as well as defaults &
  A correlation model that determines the joint probability distribution of
  rating changes for different companies, via a \textbf{Gaussian copula} \\
\bottomrule
\end{tabular}
\end{center}
\figcap{The discriminating row is the middle column. Vasicek and CreditMetrics
both model losses; CreditRisk+ models the \emph{number of defaults}, which is a
different random variable. CreditMetrics is the only one of the three that
captures downgrades short of default.}

\gapbadge

{\footnotesize\color{FRMGrey}The objective asks how each model handles default
correlation. The source notes name the Poisson and negative binomial
distributions for CreditRisk+ but do not say where either comes from or how the
model represents correlation at all. The following is written from the GARP
curriculum. CreditRisk+ was forward-referenced here from CR-5 f.}

\begin{crgapbox}[What the objective asks for: Credit Risk Plus]
In 1997 Credit Suisse Financial Products proposed a methodology for calculating
VaR that it termed Credit Risk Plus. It involves analytic approximations that are
well established in the \emph{insurance} industry.

\smallskip
Suppose an institution has $n$ loans of a certain type and the probability of
default for each loan during a year is $q$, so the expected number of defaults is
$qn$. Assuming default events are \emph{independent}, the probability of $m$
defaults is binomial. If $q$ is small and $n$ large, this is approximated by the
Poisson distribution:
\[ \text{Prob}(m \text{ defaults}) \;=\; \frac{e^{-qn}(qn)^{m}}{m!} \]
This holds approximately even when the probability of default is not the same for
each loan, provided all the default probabilities are small and $q$ is the
average probability of default over the next year for the loans in the portfolio.

\smallskip
\textbf{Where the correlation enters.} In practice the default rate $q$ is itself
uncertain and varies greatly from year to year. A convenient assumption is that
$qn$, the expected number of defaults, follows a \term{gamma distribution} with
mean $\mu$ and standard deviation $\sigma$. The Poisson distribution then becomes
a \term{negative binomial} distribution.

\smallskip
As $\sigma$ tends to zero the negative binomial converges back to the Poisson, so
the Poisson case is the zero-correlation limit. As $\sigma$ increases, the
probability of an extreme outcome involving a large number of defaults rises.
Correlation in CreditRisk+ is therefore expressed as \emph{uncertainty about the
default rate}, not as a pairwise correlation parameter.
\end{crgapbox}

% =====================================================================
\lo{CR-10 f}{Define credit spread risk and assess its impact on calculating credit VaR.}

\begin{crdefbox}[Credit spread risk]
Credit spread risk is the potential impact of changes in credit spreads on the
value of credit-sensitive products in the \term{trading book}. It affects the
probability of default and the potential loss distribution.
\end{crdefbox}

How it is incorporated depends on the method used:

\begin{itemize}
\item \term{Historical simulation.} The impact of credit spread risk can be
considered by collecting changes in credit spreads for each firm that the bank
was exposed to during the observation period.
\item \term{CreditMetrics.} A rating transitions matrix is calculated and a Monte
Carlo simulation incorporates historical data to model credit spread changes,
allowing for the calculation of portfolio value and obtaining VaR.
\end{itemize}
